% QUESTAO
Em certo jogo, cada participante lança dois dados e adiciona os números obtidos, tornando-se o ganhador aquele que obtiver o maior resultado. A tabela mostra ``alguns'' dos possíveis resultados em uma jogada.
\\\vspace{-13mm}\\
\begin{center}
\begin{tabular}{|c|c|c|c|}\hline
 & \tikz[baseline=(char.base)]{\draw[fill=gray!40] (-0.22,-0.22) rectangle (0.22,0.22);\fill (0,0) circle (0.4mm);} & \tikz[baseline=(char.base)]{\draw[fill=gray!40] (-0.22,-0.22) rectangle (0.22,0.22);\fill (-0.12,-0.12) circle (0.4mm);\fill (0.12,0.12) circle (0.4mm);} & \tikz[baseline=(char.base)]{\draw[fill=gray!40] (-0.22,-0.22) rectangle (0.22,0.22);\fill (-0.12,-0.12) circle (0.4mm);\fill (0,0) circle (0.4mm);\fill (0.12,0.12) circle (0.4mm);} \\\hline
\tikz[baseline=(char.base)]{\draw[fill=gray!40] (-0.22,-0.22) rectangle (0.22,0.22);\fill (0,0) circle (0.4mm);} & 2 & 3 & 4 \\\hline
\tikz[baseline=(char.base)]{\draw[fill=gray!40] (-0.22,-0.22) rectangle (0.22,0.22);\fill (-0.12,-0.12) circle (0.4mm);\fill (0.12,0.12) circle (0.4mm);} & 3 & 4 & 5 \\\hline
\tikz[baseline=(char.base)]{\draw[fill=gray!40] (-0.22,-0.22) rectangle (0.22,0.22);\fill (-0.12,-0.12) circle (0.4mm);\fill (0,0) circle (0.4mm);\fill (0.12,0.12) circle (0.4mm);} & 4 & 5 & 6 \\\hline
\end{tabular}
\end{center}
\\\vspace{-1mm}\\
Escreva, como sugere a tabela acima, uma matriz $R$ que represente ``todos'' os possíveis resultados desse jogo. Assim, em um jogo com $2$ participantes, se o 1º obtém $8$ como resultado, quantas possibilidades tem o 2º participante de ganhar o jogo?

% ALTERNATIVAS
$10$ possibilidades em $36$
$15$ possibilidades em $36$
$10$ possibilidades em $9$
$18$ possibilidades em $36$
$3$ possibilidades em $9$


